\section{An inertia criterion and four families with bounded exponent gaps}
\label{sec:low-spectrum}

An endpoint sign change need not detect two roots in the same interval.
For example, the complement quintic at equal exponents factors as
\[
h_{a,a,a}(x)=(x-a^2-a)\bigl(x^2-(a^2+4a)x+3a^2\bigr)^2.
\]
We instead count low eigenvalues by a symmetric Schur complement.

\begin{proposition}
\label{prop:low-spectrum}
Order the exponents as $a\leq b\leq c$, and let $m$ be an integer
with $2\leq m<a$. If
\begin{equation}
\label{eq:upper-inertia}
(c-m)(a+b+c-m)<mab,
\end{equation}
then the six-support complement quotient has exactly two nonzero
eigenvalues in $(0,m)$. If, in addition,
\begin{equation}
\label{eq:lower-inertia}
(a-m+1)(a+b+c-m+1)>(m-1)bc,
\end{equation}
both lie in $(m-1,m)$, so $G_{p^a q^b r^c}$ is not Laplacian integral.
\end{proposition}

\begin{proof}
Put $s=a+b+c$ and $P=abc$. Symmetrize $C$ by the square roots of its
cell sizes and reorder the classes as $1,2,3,23,13,12$. The resulting
real symmetric matrix is
\[
\mathcal C=\begin{pmatrix}A&-\sqrt P I\\-\sqrt P I&D\end{pmatrix},
\quad D=\operatorname{diag}(a,b,c),
\]
where $A$ has diagonal $(b+c+bc,a+c+ac,a+b+ab)$ and
off-diagonal entries $A_{ij}=-\sqrt{a_i a_j}$, with
$(a_1,a_2,a_3)=(a,b,c)$. For $0<x<a$, the block $D-xI$ is positive
definite, and Sylvester's law of inertia gives
\[
\operatorname{inertia}(\mathcal C-xI)
=\operatorname{inertia}(D-xI)+\operatorname{inertia}(K(x)),
\]
where
\[
K(x)=\operatorname{diag}(k_a(x),k_b(x),k_c(x))-vv^T,
\quad v=(\sqrt a,\sqrt b,\sqrt c)^T,
\quad k_t(x)=s-x-\frac{xP}{t(t-x)}.
\]
For fixed $x>0$, the function $k_t(x)$ increases with $t>x$, since
\[
\frac{\partial k_t(x)}{\partial t}
=\frac{xP(2t-x)}{t^2(t-x)^2}>0.
\]
Condition~\eqref{eq:upper-inertia} is exactly $k_c(m)<0$.
Thus all three diagonal values are negative, so $K(m)$ is negative
definite. The matrix $\mathcal C-mI$ has three negative and three
positive eigenvalues. The complement support graph is a triangle with
one pendant cell at each vertex; it is connected, so $C$ has a simple
zero eigenvalue and all other eigenvalues are positive. Exactly two
nonzero eigenvalues therefore lie in $(0,m)$.

For $x=m-1$, condition~\eqref{eq:lower-inertia} says $k_a(x)>0$,
so all three $k_t(x)$ are positive. Congruence by the inverse square
roots of these values turns $K(x)$ into $I-zz^T$, where
\[
\|z\|^2=\frac{a}{k_a(x)}+\frac{b}{k_b(x)}+\frac{c}{k_c(x)}>1,
\]
because each $k_t(x)<s$. Consequently $K(x)$ has one negative and two
positive eigenvalues, and the only eigenvalue of $C$ below $x$ is zero.
Both low positive eigenvalues lie in $(m-1,m)$. They are less than
$a<N$, so the complement identity and universal-vertex lift apply.
Their graph eigenvalues lie in $(V-m,V-m+1)$, which contains no integer.
\end{proof}

\begin{theorem}
\label{thm:bounded-gaps}
For every $a\geq1$, each of the exponent triples
\[
(a,a+1,a+2),\quad(a,a+1,a+3),\quad
(a,a+2,a+3),\quad(a,a+3,a+4)
\]
gives a nonintegral ideal intersection graph on three distinct primes.
\end{theorem}

\begin{proof}
For $a\leq3$, use Theorem~\ref{thm:minimum-three}. Assume $a\geq4$
and write the triple as $(a,a+d,a+e)$, where
$(d,e)=(1,2),(1,3),(2,3),(3,4)$. The left side minus the right side
of~\eqref{eq:upper-inertia}, at $m=3$, is respectively
\[
-6a,\quad-2a,\quad-4a,\quad4-2a.
\]
These values are negative. Proposition~\ref{prop:low-spectrum}
therefore places exactly two nonzero complement eigenvalues in $(0,3)$.
It remains to exclude the integers $1$ and $2$.

The first endpoints have the identities
$h_{a,a+d,a+e}(1)=-4M_{d,e}(a)P_{d,e}(a)$ from
Table~\ref{tab:gap-first-endpoints}. Every multiplier is positive for
$a\geq4$, and every coefficient of $P_{d,e}(4+u)$ is positive.
Thus $h(1)<0$ in each family.

\begin{table}[ht]
\centering
\begin{tabular}{lll}
\toprule
$(d,e)$&$M_{d,e}(a)$&Descending coefficients of $P_{d,e}(4+u)$\\
\midrule
$(1,2)$&$a(a+1)$&$(1,18,116,314,294)$\\
$(1,3)$&$a(a+1)(a+3)$&$(1,13,51,57)$\\
$(2,3)$&$a+2$&$(1,25,242,1126,2484,2022)$\\
$(3,4)$&$a+3$&$(1,28,302,1550,3699,3126)$\\
\bottomrule
\end{tabular}
\caption{Exact first-endpoint identities and positive shifted factors.}
\label{tab:gap-first-endpoints}
\end{table}

The second endpoints factor as
\begin{align*}
h_{a,a+1,a+2}(2)&=-aF_{1,2}(a),\\
h_{a,a+1,a+3}(2)&=-(a-2)(a+1)F_{1,3}(a),\\
h_{a,a+2,a+3}(2)&=-aF_{2,3}(a),\\
h_{a,a+3,a+4}(2)&=-F_{3,4}(a),
\end{align*}
where
\begin{align*}
F_{1,2}(a)&=a^5-9a^4-14a^3+93a^2+67a+6,\\
F_{1,3}(a)&=a^4-6a^3-42a^2-39a-10,\\
F_{2,3}(a)&=a^5-5a^4-48a^3-59a^2-39a-54,\\
F_{3,4}(a)&=a^6-a^5-74a^4-375a^3-967a^2-1220a-340.
\end{align*}
Table~\ref{tab:gap-modular} evaluates each factor at every residue of
one prime. No value is zero, so no factor has an integer root. The
linear multipliers are nonzero for $a\geq4$, hence $h(2)\ne0$.
Both positive eigenvalues in $(0,3)$ are therefore noninteger, and
their graph lifts are noninteger as well.

\begin{table}[ht]
\centering
\begin{tabular}{lll}
\toprule
Factor&Prime&Values at residues $0,\ldots,p-1$\\
\midrule
$F_{1,2}$&7&$(6,4,1,5,6,5,1)$\\
$F_{1,3}$&11&$(1,3,9,8,2,6,4,4,5,8,5)$\\
$F_{2,3}$&7&$(2,6,5,3,5,4,3)$\\
$F_{3,4}$&13&$(11,1,3,4,3,1,9,9,11,11,6,1,8)$\\
\bottomrule
\end{tabular}
\caption{Complete root-free modular certificates for the second endpoints.}
\label{tab:gap-modular}
\end{table}
\end{proof}

\begin{corollary}
\label{cor:span-three}
Every positive three-prime exponent triple with maximum minus minimum
at most three gives a nonintegral ideal intersection graph.
\end{corollary}

\begin{proof}
Repeated entries are covered by Theorem~\ref{thm:repeated}. After
sorting distinct entries, the only possibilities have gaps
$(d,e)=(1,2),(1,3),(2,3)$, all covered by
Theorem~\ref{thm:bounded-gaps}.
\end{proof}

\subsection{The complete finite region for minimum exponents four through seven}

The uniform cutoff makes the requested diagnostic range complete for each
of the minimum exponents $4,5,6,7$. Here the relevant quintic is
$g_B(x)=\det(xI-B)/x$ for $H$, whereas the preceding endpoint arguments
use the complement quintic $h_C(x)=-g_B(N-x)$.

\begin{corollary}
\label{cor:minimum-seven}
Every positive three-prime exponent triple with minimum entry at most
seven gives a nonintegral ideal intersection graph.
\end{corollary}

\begin{proof}
Theorem~\ref{thm:minimum-three} handles minimum entries at most three.
For a minimum entry $a\in\{4,5,6,7\}$, repeated entries are covered by
Theorem~\ref{thm:repeated}, and Proposition~\ref{prop:distinct-tail}
settles $c\geq4a^2-2a$. The remaining ordered triples lie in the exact
finite domain $a<b<c<4a^2-2a$.

For all $27\,562$ triples in this domain, the integer discriminant
$\Delta_{a,b,c}$ of $g_B$ is positive and strictly between two
consecutive integer squares. Table~\ref{tab:open-region} gives the
complete counts. The certificate
\path{results/open-region-discriminants.csv} records every triple,
discriminant and integer $q$ satisfying
$q^2<\Delta_{a,b,c}<(q+1)^2$. Each discriminant is independently
recomputed as a $9\times9$ integer Sylvester determinant by
\path{scripts/check_open_region_certificate.py}, using fraction-free
elimination and exact divisions. The certificate exhausts the domain
without omitted or duplicate triples.

If all roots of the monic quintic $g_B$ were integers, its discriminant,
the product of the squared pairwise root differences, would be an
integer square. Thus every certified quintic has a noninteger root.
The quotient is similar to a real symmetric Laplacian, so the root is
real and nonzero; Lemma~\ref{lem:join} gives a noninteger graph
eigenvalue. This complete finite certificate and the written cutoff
exhaust all triples with minimum entry at most seven.
\end{proof}

\begin{table}[ht]
\centering
\begin{tabular}{rrrrrr}
\toprule
Minimum $a$&4&5&6&7&Total\\
\midrule
Exclusive cutoff&56&90&132&182&---\\
Complete pairs&1275&3486&7750&15051&27562\\
Nonsquare discriminants&1275&3486&7750&15051&27562\\
\bottomrule
\end{tabular}
\caption{Exhaustive discriminant certificates in the region derived from
the uniform cutoff.}
\label{tab:open-region}
\end{table}

The signs of $(g_B(1),g_B(2),g_B(3))$ are $(-,-,-)$ throughout this
finite range. The complement signs instead have three patterns,
$(-,+,-),(-,+,+),(+,+,+)$, recorded separately in the diagnostic.
Different endpoint patterns do not rule out a uniform proof, and equal
signs can miss two roots in one interval. Proposition~\ref{prop:low-spectrum}
addresses that issue without extrapolating from these finite counts.
